% Inverter equivalent circuits, redrawn from Fig. 2 of the PNNL
% REGFM_A1 model specification (Du et al., PNNL-35110, 2023):
% (a) internal voltage source E<delta_E behind coupling reactance X_{\mathrm{L}};
% (b) its Norton equivalent, current source I_{\mathrm{N}}<phi_N in parallel with
% admittance Y_{\mathrm{L}}, which is what the network solver sees.
\begin{tikzpicture}[x=1cm, y=1cm, line cap=round, line join=round,
    wire/.style={line width=0.7pt},
    lbl/.style={font=\normalsize},
    >={Stealth[length=2.2mm]}]
  % ================= (a) Thevenin form =================
  \begin{scope}[shift={(0,0)}]
    % source
    \draw[wire] (0,0) circle (0.38);
    \draw[wire] plot[domain=-0.25:0.25, samples=40, smooth]
        ({\x}, {0.14*sin(deg(\x*3.1416/0.25))});
    \node[lbl, left] at (-0.45,0) {$E\angle\delta_{\mathrm{E}}$};
    % ground below source
    \draw[wire] (0,-0.38) -- (0,-0.72);
    \draw[wire] (-0.26,-0.72) -- (0.26,-0.72);
    \draw[wire] (-0.17,-0.81) -- (0.17,-0.81);
    \draw[wire] (-0.08,-0.90) -- (0.08,-0.90);
    % up and across through the reactance
    \draw[wire] (0,0.38) -- (0,1.2) -- (0.75,1.2);
    \draw[wire] (0.75,1.2)
       arc[start angle=180, end angle=0, radius=0.16]
       arc[start angle=180, end angle=0, radius=0.16]
       arc[start angle=180, end angle=0, radius=0.16]
       arc[start angle=180, end angle=0, radius=0.16];
    \node[lbl, above] at (1.39,1.38) {$X_{\mathrm{L}}$};
    \draw[wire] (2.03,1.2) -- (2.55,1.2);
    % terminal node
    \filldraw[fill=white, draw=black, line width=0.7pt] (2.55,1.2) circle (0.05);
    \node[lbl, above right] at (2.5,1.28) {$V\angle\delta_{\mathrm{V}}$};
    \draw[->, wire] (2.05,0.85) -- (2.55,0.85);
    \node[lbl, below] at (2.35,0.8) {$p,\;q,\;I\angle\varphi$};
    \node[lbl] at (1.25,-1.15) {(a)};
  \end{scope}
  % ================= (b) Norton form =================
  \begin{scope}[shift={(4.6,0)}]
    % top rail (extended so the flow arrow clears the Y_{\mathrm{L}} branch)
    \draw[wire] (0,1.2) -- (2.9,1.2);
    % current source: diamond with arrow
    \draw[wire] (0,1.2) -- (0,0.42);
    \draw[wire] (0,0.42) -- (0.30,0.05) -- (0,-0.32) -- (-0.30,0.05) -- cycle;
    \draw[->, wire] (0,-0.2) -- (0,0.3);
    \node[lbl, left] at (-0.36,0.05) {$I_{\mathrm{N}}\angle\varphi_{\mathrm{N}}$};
    \draw[wire] (0,-0.32) -- (0,-0.72);
    % admittance: rectangle
    \draw[wire] (1.5,1.2) -- (1.5,0.5);
    \draw[wire] (1.32,0.5) rectangle (1.68,-0.4);
    \node[lbl, right] at (1.72,0.05) {$Y_{\mathrm{L}}$};
    \draw[wire] (1.5,-0.4) -- (1.5,-0.72);
    % bottom rail + ground
    \draw[wire] (0,-0.72) -- (1.5,-0.72);
    \draw[wire] (0.75,-0.72) -- (0.75,-0.86);
    \draw[wire] (0.49,-0.86) -- (1.01,-0.86);
    \draw[wire] (0.58,-0.95) -- (0.92,-0.95);
    \draw[wire] (0.67,-1.04) -- (0.83,-1.04);
    % terminal
    \filldraw[fill=white, draw=black, line width=0.7pt] (2.9,1.2) circle (0.05);
    \node[lbl, above right] at (2.85,1.28) {$V\angle\delta_{\mathrm{V}}$};
    \draw[->, wire] (2.4,0.85) -- (2.9,0.85);
    \node[lbl, below] at (2.7,0.8) {$p,\;q,\;I\angle\varphi$};
    \node[lbl] at (1.0,-1.35) {(b)};
  \end{scope}
\end{tikzpicture}
